\chapter{Corpos}

\noindent{}O filósofo e historiador Martin Buber conta que certa vez perguntaram ao
Rabi Iehiel Mihal\footnote{Rabi Iehiel Mihal era conhecido como
  \emph{Tzadik HaGamur} (um Justo) e foi um dos mais notáveis discípulos
  do Rabi Baal Shem Tov, precursor do movimento \emph{hassídico} judaico
  no Leste Europeu.} a razão de retardar o começo de suas orações. O
rabino respondeu: ``Como a tribo de Dan\footnote{Após o
  recebimento das Tábuas da Lei, o povo hebreu foi dividido em doze
  tribos, sendo cada uma delas responsável por um pedaço de terra em
  \emph{Canaã} (antigo território que hoje engloba Israel, Palestina,
  Líbano, Cisjordânia, Jordânia e o Deserto do Sinai). As doze tribos
  são: Reuven, Shimon, Levi, Yehuda, Dan, Naftali, Gad, Asher,
  Issachar, Zebulun, Yossef e Biniamyn.}, que na Terra de Canaã
caminhava atrás de todas as demais tribos, recolhendo tudo aquilo que
perdiam; assim também sou eu --- um coletor incansável de migalhas de
nossa, e minha, própria história''\footnote{\textsc{buber}, \emph{op.\,cit.}, p.\,192.}.

Tal qual a tribo de Dan, Rabi Iehiel Mihal e Martin Buber, no judaísmo
todos os indivíduos são convidados a fazer um \emph{chidush}\footnote{Na
  literatura rabínica, o \emph{chidush} refere-se a uma nova abordagem
  de ideias e interpretações em obras já existentes. É uma forma de
  inovação feita dentro do sistema da \emph{Halachá} --- série de
  documentos que descrevem as maneiras judaicas de conduzir a vida, suas
  regras e verdades.} e podem, assim, se apropriar e interpretar os
textos sagrados e bíblicos\footnote{Na tradição judaica todos os
  indivíduos possuem o mesmo nível de conexão com a divindade, sem a
  necessidade de intermediação.}. Para se tornar este coletor de
migalhas --- vestígios de uma história particular compartilhada ---,
como nos lembram o escritor Amós Oz e a historiadora Fania
Oz-Salzberger, não é preciso ser filha ou filho de um ventre judaico,
basta ser um leitor --- um apropriador incansável, um intenso
talmúdico, um intérprete.

Das mais dispersas migalhas de narrativas, individuais e coletivas,
agarradas em meio ao trânsito, essas múltiplas referências, por fim, se
tornam sementes e são germinadas em determinadas territorialidades,
agregando influências e interferências dos locais que passam a
frequentar. Pois, no campo dos Estudos Judaicos, entende-se que
historiadoras e historiadores não produzem uma História Judaica por
excelência, mas sim histórias --- narradas e escritas --- sobre as
pegadas destes judeus e judias nos territórios que percorrem.

Neste sentido, este texto sustenta-se enquanto um novo começo, de outras
--- mas nunca novas --- interpretações passíveis de apropriação e de
pertencimento. Como todas as evidências são escoradas por mistérios,
entende-se aqui que se não se fizer o fogo não há o que ver, ou melhor,
não há de se ver nada que se deseja. Porém, ao se decidir por acendê-lo
então \emph{kudiado con hacer fuego, i no aguentar el kalor}.

\begin{center}\adforn{49}\end{center}

\noindent{}Em ``Tempo e atemporalidade'', terceiro capítulo do livro \emph{Os
judeus e as palavras}, algumas dessas ``clareiras-fogos'' já se mostram
acesas. Amós Oz e Fania Oz-Salzberger discorrem a respeito do conceito
de espaço e de tempo para a civilização, para a tradição e para a
cultura judaicas. Parte-se do pressuposto de que todas as civilizações
estão preocupadas com o seu passado e no judaísmo isso não poderia ser
diferente. São as orações e as rezas, os cultos, os ritos, as tradições,
a cultura e as festividades que lembram e relembram cotidianamente o
contínuo \emph{quórum} --- em que a prática levaria à mística\footnote{Os
  ritos judaicos baseiam-se em uma série de preceitos e de tarefas
  religiosas realizadas cotidianamente. A \emph{kavaná} --- a elevação
  mental e espiritual que o indivíduo atinge ao executar os deveres
  religiosos --- está também entremeada de uma rotina, em que a
  constante prática e o refazimento são essenciais para alcançar a
  santidade.} --- de uma tradição ancorada em textos e em livros.

Assim, em parte, é o fio genealógico --- gênico ou textual --- que
transmite as mais variadas interpretações e tradições judaicas, de
geração em geração --- e firmam-se, finalmente, para que então possam
ser lidas e tateadas em outros territórios e em outros
tempos\footnote{Aqui vale lembrar que os documentos, as tradições e as
  comunidades, ao firmarem-se em outras territorialidades e tempos,
  ainda são passíveis de indeterminabilidade, condicionados aos
  movimentos das diásporas voluntárias ou forçadas. Ao serem levados
  novamente pelo velho vento do sopro nômade da diáspora, brotarão em
  outros locais, em outros tempos, em um movimento interminável, cíclico
  e constante.}. Em ressonância, todos os dias da semana, ao colocar o
\emph{tefilin}\footnote{O \emph{tefilin} é um objeto de cunho religioso
  que consiste em um par de caixas unidas por meio de uma tira de couro
  animal. Conhecido como filactério, o \emph{tefilin} é um instrumento
  de proteção e fortificação da aliança divina usado por homens, e mais
  recentemente incorporado por mulheres, nos momentos de culto e rito
  religioso. Esse objeto é utilizado apenas durante os dias de semana,
  não sendo incorporado nas práticas judaicas do \emph{shabat}.}, judeus
e judias selam a devoção e afeição entre Deus e o povo judeu. Na
Torá, Deus e povo são aludidos enquanto um casal, um par --- e,
em união, simbolizam a aliança, a lealdade intelectual e uma
espiritualidade e religiosidade atreladas ao divino\footnote{Segundo o
  Talmud, Deus criou os seres humanos --- e todas as coisas do
  mundo --- para terem prazer com a conexão divina. Todas as pessoas são
  passíveis de manifestar o divino em suas vidas por meio de ações
  ordinárias e extraordinárias. Pois entende-se que o divino --- para
  além de Deus --- é a presença onipotente e onisciente, e é ele,
  propriamente, a matéria.}.

É a partir do conceito de que a divindade está permeada de livros, de
textos e de ritos --- em especial uma lealdade intelectual --- que o
Rabi Akiva\footnote{Rabi Akiva (50 d.E.C.--135 d.E.C.) foi um renomado
  estudioso hebreu e introduziu um novo método de interpretação das leis
  e das tradições orais judaicas, que veio a se tornar a \emph{Mishná}. A \emph{Mishná} é uma das principais obras do judaísmo rabínico
  e a primeira grande redação da tradição oral judaica.}, ao se deparar
com a destruição do segundo Templo de Jerusalém\footnote{O segundo
  Templo Sagrado de Jerusalém, ou Templo de Jerusalém, foi uma série de
  construções de cunho religioso que se localizava no Monte
Moriá (cidade velha de Jerusalém), atual local do Domo da
  Rocha e da Mesquita \emph{Al-Aqsa}.} (76 d.E.C.), dança incansável e
insaciadamente em torno de seus escombros e resíduos. É instigante
pensar que, das duas vezes\footnote{O Templo de Jerusalém foi destruído
  duas vezes, no ano de 587 a.E.C. e posteriormente no ano de 76 d.E.C.}
que se tentou --- talvez a todo o custo --- fixar a
\emph{Shechiná}\footnote{Conta o Talmud que, quando os judeus
  pecaram, a presença de Deus deixou o Templo Sagrado de Jerusalém e
  retirou-se em direção aos céus. Porém, uma parte da \emph{Shechiná}
  permaneceu em terra para acompanhar os judeus e judias que
  forçosamente deixaram a Terra de \emph{Canaã}, levados ao exílio na
  Babilônia. Conta a tradição que a presença divina --- não mais
  presidida em uma construção de pedra --- passou a ser transportada
  pelas crianças¹ em direção ao exílio. Outro exemplo diz: ``Está
  escrito na Guemará: Todas as datas previstas já passaram. Mas assim
  como a Shechiná deixou o Santuário e foi ao exílio no decurso de dez
  viagens, assim também ela não voltará de uma só vez e a luz da
  salvação paira entre o céu e a terra''. A \emph{Shechiná} foi levada ao exílio através das crianças,
  que, na tradição judaica, são instrumentos visíveis de esperança, de
  que haverá uma continuidade deste povo. \textsc{buber}, \emph{op.\,cit.}, p.\,386.}, a habitação e morada de Deus foi
destruída. Assim, a divindade que sobrevive a esta destruição --- e que
neste texto busca-se reconstituir --- está diretamente vinculada a um
conceito de Deus e de divino amparado por arquiteturas não mais estáveis
e intransponíveis, mas sim móveis, reprogramadas nos espaços e nos
tempos, apropriadas por cada uma das gerações judaicas e levadas em
bolsos, maletas, sacolas nas diásporas.

Ainda nesta chave, como nos lembram Amós Oz e Fania Oz-Salzberger:
Não existe \emph{terra santa} na Torá. O único lugar sacrossanto é o
Templo de Jerusalém''\footnote{\textsc{oz; oz-salzberger}, \emph{op.\,cit.}, p.\,121.}
--- motivo pelo qual todas as sinagogas estão voltadas nessa direção. A
partir das constantes destruições --- de arquiteturas e de comunidades
---, o renascer das cinzas das ruínas faz parte das vidas judaicas em
diásporas e se estabelece enquanto pedras --- estilhaçadas e
fragmentadas --- com tanta sede de celebração que são capazes de
inventar a vida só para poder beber dessa existência. 

E se essas pedras falassem, o que elas diriam?\footnote{Elas diriam que:
``No verdadeiro amor tudo é para sempre vivo. E sei que, como as
  pedras, existo pela sede. Quero sempre inventar a vida''. \textsc{mãe}, \emph{As mais belas coisas do mundo}, p.\,43.} Antes mesmo
de selar um contrato matrimonial --- um dos principais motivos de
celebração das comunidades judaicas ortodoxas --- há o costume de o
noivo quebrar um copo de vidro contra o chão. Este ato é uma recordação
da destruição do Templo Sagrado e de que nunca se estará completo sem a
presença da \emph{Shechiná}\footnote{Conta Martin Buber: ``Quando
  o Rabi Baruch chegou, no salmo, às palavras: \emph{Não deixarei que durmam
  meus olhos, nem que descansem minhas pálpebras antes de ter encontrado
  uma morada para Deus}, parou e disse de si para consigo: \emph{Antes de eu
  me ter encontrado e transformado em morada pronta para o pouso da
  Shechiná}''. \textsc{buber}, \emph{op.\,cit.}, p.\,131.}. Por fim, quebra-se o copo e dança-se
em seus escombros, como fez o Rabi Akiva.

O judaísmo é feito de tempos santos. Eles se consolidam enquanto uma
espécie de calendário cíclico e elipsoidal que foge do binômio
progresso-tempo, abruptamente implementado pelo renascimento e pela
modernidade ocidental. Entre estes tempos santos estão o \emph{shabat},
os \emph{chaguim}, os jejuns, os dias de \emph{Yom Tov}, as mais
diversas e variadas celebrações e festividades cíclicas judaicas --- que
recordam desastres e perseguições, momentos de glórias e triunfos,
salvamentos divinos e outros salvamentos essencialmente humanos. Ao fim,
entende-se que, embora o Templo de Jerusalém tenha sido destruído, os
judeus ``ainda puderam levar consigo para a amarga diáspora suas
santidades não espaciais, intangíveis: a língua, as leituras e o sempre
recorrente, ciclicamente confortante calendário de
\emph{tempos-santos}''\footnote{\textsc{oz; oz-salzberger}, \emph{op.\,cit.}, p.\,121.}.

Estes tempos santos --- alguns recordados semanalmente, outros
anualmente\footnote{Alguns tempos sagrados judaicos são lembrados e
  realizados semanalmente, como o \emph{shabat} --- recordação da
  criação divina do mundo e de todas as coisas que ele contém. Outros
  eventos são recordados anualmente, como as festividades e os jejuns.
  Há ainda as recordações geracionais, que ocorrem a cada sete anos,
  como é o caso dos anos de \emph{shemitá}, em que há a proibição do
  trabalho agrário, permitindo-se assim apenas a coleta dos frutos e
  frutas.} --- são celebrações de rito e de vida partilhadas e
realizadas em comunidade. Ao final, entende-se fortemente que os povos
em diáspora celebram a vida. E, como uma espécie de calendário,
espera-se que, um dia, no lugar de memórias difíceis e tristeza --- tais
quais os eventos de perseguição, de intolerância, de destruição ---
haverá celebração. Mas não esperaremos calados.

O artifício tão recorrentemente usado de celebrar a vida\footnote{O
  humor é uma característica intrínseca aos judeus \emph{ashkenazi}
  (judeus provenientes da Europa Central e do Leste Europeu). Apesar de
  ser uma particularidade relativamente recente e incorporada à cultura
  judaica no final do século \textsc{xix}, este atributo foi introduzido na
  literatura iídiche por meio do escritor Sholem Aleichem ---
  que, ao retratar o sofrimento do povo judeu, se viu forçado a manter
  acesa a chama da vida, da alegria e do humor de viver. Assim:
  ``{[}\ldots{}{]} Sholem Aleichem capturou a alma de um povo que se
  viu identificado e que sorriu em retorno para ele e aprendeu a sorrir
  com ele''. \textsc{kilsztajn}, \emph{op.\,cit.}, p.\,79.} --- um \emph{modus operandi}
judaico de ser e estar no mundo --- está presente nos festejos ortodoxos
e tradicionais, nas banalidades do cotidiano, na quebra do copo de vidro
ao selar um contrato matrimonial. Está firmado no sempre alto, e sempre
belo, \emph{Lechaim}\footnote{A palavra hebraica \emph{Lechaim}, usada
  nos momentos de celebração, é escrita e falada no plural (todas as
  palavras hebraicas terminadas com sufixos ``im'' ou ``ot'' são
  palavras no plural). Pois, no judaísmo, entende-se que o
  \emph{Lechaim} (no português: ``Saúde!'') deve sempre ser entoado
  coletivamente, em comunhão. Uma segunda interpretação diz que, quando
  celebra-se a vida, não se deve celebrar somente uma única vida. Todas
  as vidas são passíveis de celebração e de cerimônia. Dessa maneira, a
  língua hebraica compreende que o \emph{Lechaim}, melhor traduzido para
  ``Às vidas!'', para além de uma celebração compartilhada entre pessoas
  que situam-se no mesmo espaço-tempo, é uma forma de festejar e lembrar
  todas as demais existências criadas e moldadas por Deus.}. Como também
nas orações, nas bênçãos e nos cânticos entoados coletivamente --- em
comunhão --- nas mais dispersas e variadas sinagogas e comunidades. O
celebrar está nos ritos, na passagem de morte. São temporadas curtas,
mas alegres, que se repetem anualmente, e que são guardadas e
transmitidas pelos movimentos cíclicos e elipsoidais da continuidade de
uma tradição, e de uma educação judaica, que garante constantes
manutenções --- que atravessa escalas, e é ao mesmo tempo transversal a
elas --- de vínculos pessoais de identidade, que consolidam comunidades
nos mais diversificados territórios.

Dessa maneira, ao traçar este raciocínio, o \emph{chidush}, a
continuidade textual, a lealdade intelectual, o senso de comunidade, de
agregação, de coletivo, as leituras e os estudos bíblicos e
talmúdicos são preceitos que garantem vínculos identitários e
comunitários entre as diversas gerações que transitam pelo mesmo
espaço-tempo dentro das comunidades. Assim, esses sujeitos --- ao se
debruçarem sobre os textos sagrados --- transcendem as barreiras
temporais e criam para si outros mundos --- outros lugares de diálogo e
de estudos por meio das interpretações.

% \begin{center}\adforn{49}\end{center}